
\section{Project Objectives and Requirements}

The goal of the project is to define heat schedules for all available production units with the lowest possible expenses and the highest profit on the electricity market. For the development, the SCRUM framework was adopted. Communication between the team and product owners was established early on to provide a better feedback loop and an opportunity for thorough requirement elicitation, which were compiled in a table (Appendix \ref{appendix:requirements}), and those prescribed in two scenarios. The user should be able to see visual representation of all relevant data, assets, optimization results, as well as being able to configure each individual asset, optimizer itself and export the results in a convenient format. The system shall ensure that heat demand is met at any time, calculate net cost for each production unit and prioritize each unit based on it, create separate result data for each production unit and calculate the optimized heat schedule. 

\section{Problem Statement}

HeatItOn utility company is supplying heat to a single district heating network with nearly two thousand buildings. They utilize different units for heat production, such as traditional gas boilers or units that either consume or produce electricity. At the moment, the company does all its scheduling of units manually. However, manual scheduling is both time consuming and prone to human error. Managing all the units, constantly following the electricity prices that change hourly, and ensuring that the heat demand is met for all the buildings on the grid makes it harder to optimize the production for the best profit. Moreover, the time spent on this task could be better utilized on other, higher-value activities.\par
An automated system for optimizing the production for profit can significantly speed up the scheduling of units and save human resources otherwise spent on manual calculations. This would ensure that the heat production is as profitable as it can be and improve accuracy by reducing human error. The project aims to develop a heat optimizer software that can take in all the units and their data and calculate a unit schedule that will be optimized for the highest profit. It takes into consideration the hourly shifting electricity prices and the different heating seasons to make the best and most profitable choices for the unit schedule.

